% paper/sections/02_methods.tex
% Methods for the JCTC revision of ct-2026-01838j.
% Derived numbers are \fs... macros from paper/macros.tex. Training recipes,
% dataset sizes and protocol settings are transcribed from
% paper/REVISION_FACTS.md (sections 4 and 7) and protocols/j10_ood.yaml,
% protocols/j11_active_learning.yaml. \fsTODO marks a number that will exist as
% a macro once the J10 scoring record is committed.

\section{Methods}
\label{sec:methods}

This section defines every object the Results use.
Section~\ref{sec:committees} describes how each committee was built.
Section~\ref{sec:scores} defines the uncertainty of a committee and the three
ways we reduce that uncertainty to one number per structure.
Section~\ref{sec:baselines} defines two zero-cost baselines that every
committee score is compared against. Section~\ref{sec:metric} defines what
counts as an error and how a score is judged for ranking errors, and
Section~\ref{sec:calibmetrics} defines the calibration metrics.
Section~\ref{sec:screening} defines the cheap estimators and the screening rule
built on them. Section~\ref{sec:data} lists the datasets, the splits, and the
overlap between training and evaluation sets. Section~\ref{sec:protocols}
gives the two experiments that were pre-registered for this revision.
Table~\ref{tab:notation} collects the notation.

\subsection{Committee models}
\label{sec:committees}

A committee is a set of $M$ force predictors evaluated on the same structure.
The spread of the $M$ predictions is the uncertainty. We use three kinds of
committee, all with $M=8$ members. The first two kinds were trained by
\citet{beck2025committee} and released with the supporting data of that paper;
we use those checkpoints unchanged, and the recipes below are taken from the
training scripts and configuration files in that archive. The third kind places
new committee heads on a frozen foundation-model trunk.

\subsubsection{Shared-trunk committees}

A shared-trunk committee, also called a multi-head committee, attaches $M$
readout heads to one message-passing trunk~\citep{beck2025committee}. The
construction is a shallow ensemble in the sense of
\citet{kellner2024shallow}: the heads share every feature and differ only in
the last layer. The trunk is a MACE model~\citep{batatia2022mace} with two
interaction layers, 32 channels, maximum rotational order $L=1$, correlation
order 3 and a cutoff radius of $6.0$~\AA. Each head $m$ is a separate readout
attached to the atomic features of every layer: a linear readout on the first
layer and a multilayer perceptron with one hidden layer of 16 SiLU units on the
second. Head $m$ sums its per-atom readouts into an energy $E_m(x)$ for a
structure $x$, and the forces of head $m$ are
$F_m(x)=-\nabla_x E_m(x)$. One evaluation of the trunk therefore gives $M$
energies and, after the backward passes described in
Section~\ref{sec:screening}, $M$ force vectors.

Every head is trained with the loss
\begin{equation}
\mathcal{L}=10\,\mathrm{MSE}(E)+1000\,\mathrm{MSE}(F),
\label{eq:loss}
\end{equation}
with the per-atom normalisation of the energy term used by MACE, on the
structures assigned to that head. Gradients from every head update the trunk.
Training used the Adam optimiser with an exponential moving average of the
weights, batch size 10, float32 arithmetic and random seed 123, followed by
stochastic weight averaging (SWA) over the final quarter of the epochs.

The heads differ only through the training data each head sees. The trunk
sees every training structure in all three constructions.
\begin{itemize}
\item \emph{Disjoint heads.} The training set is divided evenly among the
  eight heads with no structure shared. Each structure updates the trunk and
  the one head assigned to that structure, so each head is fitted to one
  eighth of the data.
\item \emph{Overlapping heads.} Each head is trained on an independently drawn
  random 80\% of the training set, with the remaining 20\% used as the
  validation set of that head. Any two heads share most of their data.
\item \emph{Identical data.} Every head is trained on the full training set.
  The heads differ only by their random initialisation.
\end{itemize}
The number of epochs, with the epoch at which SWA began in parentheses, was:
3BPA disjoint 500 (375), overlapping 78 (59), identical data 65 (49); liquid
water disjoint 200 (150), overlapping 31 (23); rMD17 disjoint 320 (240). The
overlapping and identical-data committees were trained for fewer epochs because
each head sees more data per epoch.

\subsubsection{Independently trained committees}

An independently trained committee consists of eight complete MACE models of
the architecture above, each trained from scratch with the loss of
Eq.~(\ref{eq:loss}) and the same optimiser, from random seeds 0 to 7
\citep{krogh1995neural,lakshminarayanan2017simple}.
\citet{beck2025committee} call this construction a naive committee. Two
variants match two of the shared-trunk constructions. In the identical-data
variant all eight models use the same training and validation files, which
matches identical-data heads. In the resampled variant each model draws its
own random 80\% of the training set, which matches overlapping heads. The
number of epochs (SWA start) was 3BPA resampled 625 (469), 3BPA identical data
521 (391), water resampled 250 (188). The committee mean, the head variance and
every score in Section~\ref{sec:scores} are computed from the eight force
vectors by the same code as for the shared-trunk committees, so construction
is the only variable in a matched comparison.

\subsubsection{Foundation committees}

The third kind places committee heads on a foundation model. The trunk is
MACE-MP-0b (medium, agnesi radial basis)~\citep{batatia2023foundation}, which
was pretrained on the MPtraj dataset~\citep{deng2023chgnet}. The trunk is
frozen. Eight new disjoint readout heads, about $15{,}000$ parameters in
total, are the only trainable weights. Fine-tuning ran for 1000 epochs with SWA
from epoch 800, learning rate 0.01, AMSGrad, an exponential moving average with
decay 0.99, float64 arithmetic, and the loss
$10\,\mathrm{MSE}(E)+1000\,\mathrm{MSE}(F)+1\,\mathrm{MSE}(\text{stress})$,
with the isolated-atom reference energies of the foundation model. Four
variants, all from \citet{beck2025committee}, differ only in which MPtraj
structures the heads were fitted on: $8{,}000$ structures selected by
query by committee (QbC-8000), $1{,}000$ selected the same way (QbC-1000),
$8{,}000$ drawn at random (random-8000), and the $8{,}000$ with the largest
predicted mean force (max-force-8000). Because the trunk was pretrained on
MPtraj, no MPtraj structure is new to the trunk. Section~\ref{sec:data} states
what that overlap means for each experiment.

\subsection{Uncertainty and the three structure-level scores}
\label{sec:scores}

For a structure $x$ with $N$ atoms, let $F(x)\in\mathbb{R}^{3N\times M}$ hold
the predicted forces, one column per head, and let
$\bar f(x)=M^{-1}\sum_{m}F_{\cdot m}(x)$ be the committee mean. The committee
mean is the force a simulation would use. Index the $3N$ Cartesian coordinates
by $d$ and the atoms by $a$. The per-coordinate head variance is the unbiased
variance over heads,
\begin{equation}
v_d(x)=\frac{1}{M-1}\sum_{m=1}^{M}\bigl(F_{dm}(x)-\bar f_d(x)\bigr)^2 ,
\label{eq:vd}
\end{equation}
and the per-atom standard deviation is the mean over the three Cartesian
components of atom $a$,
\begin{equation}
\sigma_a(x)=\frac{1}{3}\sum_{d\in a}\sqrt{v_d(x)} ,
\label{eq:sigma}
\end{equation}
which is the per-atom uncertainty of \citet{beck2025committee}. Every
per-atom uncertainty in this paper is $\sigma_a$, and every per-coordinate
uncertainty is $\sigma_d=\sqrt{v_d}$.

A selection rule needs one number per structure. We compare three reductions
of the head variance to such a number:
\begin{align}
S_{\mathrm{glob}}(x)&=\sum_{d=1}^{3N} v_d(x) \qquad\text{(global score),}
\label{eq:global}\\
S_{\mathrm{comp}}(x)&=\max_{d}\sqrt{v_d(x)} \qquad\text{(max-component score),}
\label{eq:maxcomp}\\
S_{\mathrm{atom}}(x)&=\max_{a}\sigma_a(x) \qquad\text{(max-atom score).}
\label{eq:maxatom}
\end{align}
The global score is the sum over all $3N$ coordinates of the per-coordinate
head variance. The max-component score is the largest per-coordinate standard
deviation, and is the rule \citet{beck2025committee} used for acquisition. The
max-atom score is the largest per-atom standard deviation. The global score
sums over coordinates and therefore grows with $N$ even when the uncertainty
per coordinate is constant; the two maximum scores do not. For pools in which
$N$ varies we also use the size-normalised score
\begin{equation}
S_{\alpha}(x)=\frac{S_{\mathrm{glob}}(x)}{(3N)^{\alpha}},
\label{eq:salpha}
\end{equation}
with the exponent $\alpha$ chosen on a design split
(Section~\ref{sec:data}). The value $\alpha=0$ recovers the global score and
$\alpha=1$ gives the mean per-coordinate variance. At fixed $N$ every $\alpha$
gives the same ranking, so the normalisation changes nothing on a single-size
dataset.

The global score depends on $F(x)$ only through a small matrix. Let
$Q_c\in\mathbb{R}^{M\times r}$ with $r=M-1$ be an orthonormal basis of the
subspace of head space orthogonal to the all-ones vector; we use the Helmert
basis. Multiplying by $Q_c$ removes the mean force, and
\begin{equation}
A(x)=Q_c^{\top}F(x)^{\top}F(x)\,Q_c\in\mathbb{R}^{r\times r},\qquad
v_d(x)=\frac{\lVert Q_c^{\top}F_{d\cdot}(x)^{\top}\rVert^{2}}{M-1},\qquad
S_{\mathrm{glob}}(x)=\frac{\operatorname{tr}A(x)}{M-1},
\label{eq:A}
\end{equation}
where $F_{d\cdot}(x)\in\mathbb{R}^{1\times M}$ is row $d$ of $F(x)$. This
centred head-space matrix is needed because every cheap estimator of the
global score in Section~\ref{sec:screening} is a function of $A(x)$ alone:
$A(x)$ has $r\times r$ entries whatever the value of $N$, so a pool scan can
store $A(x)$ for every structure instead of $F(x)$, and the exactness of each
estimator can be checked against Eq.~(\ref{eq:vd}) directly. The two maximum
scores are not functions of $A(x)$ and are computed from $F(x)$.

\subsection{Zero-cost baselines}
\label{sec:baselines}

A committee score should not be compared with random selection alone, because
two informative quantities are available before any committee uncertainty has
been computed.

The \emph{predicted force norm} is the largest mean-force magnitude on any
atom,
\begin{equation}
b_F(x)=\max_{a}\lVert \bar f_a(x)\rVert .
\label{eq:bF}
\end{equation}
The mean force is the force the integrator needs, so $b_F$ costs nothing beyond
the simulation step itself. Force errors tend to grow with force magnitude, so
$b_F$ is a plausible predictor of error, and selection by the largest mean
force is one of the acquisition baselines of \citet{beck2025committee}.

The \emph{head-energy spread} is the unbiased standard deviation of the $M$
head energies,
\begin{equation}
b_E(x)=\Bigl(\frac{1}{M-1}\sum_{m=1}^{M}\bigl(E_m(x)-\bar E(x)\bigr)^2\Bigr)^{1/2}.
\label{eq:bE}
\end{equation}
The head energies are produced by the forward pass, so $b_E$ costs one forward
pass and no backward pass beyond the one that produces the mean force.

Both quantities are zero-cost relative to committee force uncertainty, which
for a shared-trunk committee needs $r=M-1$ additional backward passes
(Section~\ref{sec:screening}) and for an independently trained committee needs
$M$ complete model evaluations.

\subsection{Error targets and the ranking metric}
\label{sec:metric}

Let $F_{\mathrm{ref}}(x)$ be the reference forces from the electronic-structure
method that labelled the dataset, and let
\begin{equation}
\Delta(x)=\bar f(x)-F_{\mathrm{ref}}(x)\in\mathbb{R}^{3N}
\label{eq:delta}
\end{equation}
be the error of the committee mean. We measure the error of the committee
mean, not of the individual heads, because the mean is the prediction a
simulation uses and the prediction a selection rule is meant to protect. The
error vector is reduced to one number per structure in four ways, which we call
the error targets:
\begin{align}
e_{\max}(x)&=\max_{a}\lVert\Delta_a(x)\rVert \qquad\text{(largest per-atom error; primary),}\\
e_{\mathrm{rmse}}(x)&=\Bigl(\frac{1}{3N}\sum_{d}\Delta_d(x)^2\Bigr)^{1/2} \qquad\text{(root-mean-square error over coordinates),}\\
e_{\mathrm{maxcomp}}(x)&=\max_{d}\lvert\Delta_d(x)\rvert \qquad\text{(largest per-coordinate error),}\\
e_{\mathrm{q95}}(x)&=\text{the 95th percentile of }\lvert\Delta_d(x)\rvert\text{ over the $3N$ coordinates.}
\end{align}
The largest per-atom error is primary because a simulation fails through its
worst atom. The other three targets test whether a conclusion depends on the
reduction chosen.

For a given error target, the \emph{positives} of an evaluation set are the
5\% of structures with the largest value of that target. A score is judged by
the area under the receiver operating characteristic curve (AUROC): the
probability that a randomly chosen positive structure receives a higher score
than a randomly chosen negative structure, computed as the Mann--Whitney
statistic with tied ranks averaged. A value of 0.5 means the score carries no
information about which structures are positive; a value of 1 means every
positive is scored above every negative. AUROC depends only on the ranking, so
it is unaffected by any monotone rescaling of a score, and in particular by
the magnitude calibration of the committee. A \emph{random null}, a uniformly
random score with a recorded seed, is reported next to every comparison; its
departure from 0.5 shows the finite-sample noise of the evaluation set.

The ranking comparisons fit nothing, so they use the full evaluation set of
each system. Confidence intervals come from a moving-block
bootstrap~\citep{kunsch1989jackknife}, with block length $\lceil 2\tau_{\mathrm{int}}\rceil$
frames, where $\tau_{\mathrm{int}}$ is the integrated autocorrelation time of
the error target along the trajectory (Section~\ref{sec:data}). The bootstrap
is paired: the same resample of structures is scored with both signals of a
difference, so the interval of the difference accounts for the correlation
between the two signals. The number of resamples and the seed are stored with
each record. A difference is called significant when its 95\% interval excludes
zero. These intervals quantify the resampling of structures at a fixed split
assignment; they do not include the variability of redrawing the split.

\subsection{Calibration metrics}
\label{sec:calibmetrics}

Ranking quality and calibration are different
properties~\citep{pernot2022long,tran2020methods}: a committee can order
structures correctly while understating every error by the same factor.
Table~\ref{tab:accuracy} reports, for every committee and evaluation set, the
force mean absolute error (MAE) and root-mean-square error (RMSE) of the
committee mean and the per-atom energy MAE and RMSE, next to the following four
calibration measures. Figure~\ref{fig:calibration} shows the
error-versus-uncertainty relation that the measures summarise.

\emph{Overconfidence ratio.} The force RMSE of the committee mean over all
coordinates and structures of an evaluation set, divided by the median of
$\sigma_a$ over all atoms and structures of that set. A ratio of 1 means the
typical disagreement matches the typical error; a ratio above 1 means the
committee understates its error by that factor.

\emph{Z-score variance.} With $z_d=\Delta_d/\sigma_d$ per coordinate, we report
the mean of $z_d^2$ over all coordinates and structures. If the errors were
Gaussian with standard deviation $\sigma_d$ this variance would be 1; a larger
value means overconfidence and a smaller value
underconfidence~\citep{kahle2022quality,busk2021calibrated}. The per-atom
version uses $\lVert\Delta_a\rVert/(\sqrt{3}\,\sigma_a)$.

\emph{Spearman rank correlation.} Per atom, the Spearman correlation between
$\sigma_a$ and $\lVert\Delta_a\rVert$ pooled over all atoms of all structures.
Per structure, the Spearman correlation between $\max_a\sigma_a$ and
$e_{\max}$. A rank correlation is used for the same reason as AUROC: it is
unaffected by the magnitude calibration.

\emph{Reliability curve.} All atoms of an evaluation set are sorted by
$\sigma_a$ and divided into ten bins with equal numbers of atoms. For each bin
we plot the mean of $\sigma_a$ against the root mean square of
$\lVert\Delta_a\rVert$~\citep{kahle2022quality,levi2022evaluating}. For a
calibrated committee with Gaussian errors the root mean square of the
three-component error norm equals $\sqrt{3}\,\sigma_a$, and that line is drawn
as the reference. Equal-count bins avoid sparsely populated bins at large
$\sigma_a$, where a curve built from equal-width bins is dominated by a few
atoms.

\subsection{Cheap estimators and the screening rule}
\label{sec:screening}

\subsubsection{Cost model}

For a shared-trunk committee the forces of all heads come from one trunk, and
the cost lies in the backward passes. One reverse-mode pass through the trunk,
with the $M$ head energies weighted by a vector $w\in\mathbb{R}^{M}$, returns
the projection
\begin{equation}
g=F(x)\,w\in\mathbb{R}^{3N}.
\label{eq:projection}
\end{equation}
The mean force is the projection with $w=\mathbf{1}/M$ and costs one backward
pass, which the simulation pays in any case. The exact per-coordinate variances
of Eq.~(\ref{eq:vd}) need the $r=M-1$ projections onto the columns of $Q_c$,
hence $r$ further backward passes, after which $A(x)$ follows exactly. A cheap
estimator uses $K<r$ projections onto directions $w_1,\dots,w_K$ in the centred
subspace. Throughout, $K$ denotes the number of head-space projections an
estimator uses, and the total cost of the cheap score is $K+1$ backward passes
including the mean force. Estimators are compared at equal numbers of backward
passes. For an independently trained committee there is no shared trunk;
obtaining $F(x)\,w$ needs all $M$ columns of $F(x)$, so no estimator can cost
less than the exact quantity and the screening rule below does not apply.

\subsubsection{Estimators of the global score}

Every estimator below is a function of $A(x)$ (Eq.~\ref{eq:A}), because
$\lVert g_k\rVert^2=w_k^{\top}F(x)^{\top}F(x)\,w_k$. We compare four.
\begin{enumerate}
\item \emph{Random sketch.} In the Haar sketch, $w_1,\dots,w_K$ are an
  orthonormal set drawn uniformly at random from the centred subspace, and the
  estimate is $(r/K)\sum_{k}\lVert g_k\rVert^2/(M-1)$, which is unbiased for
  $S_{\mathrm{glob}}$. In the Gaussian sketch the $w_k$ are independent
  standard Gaussian vectors in the centred subspace with the matching scaling.
\item \emph{Head subsampling.} $K$ of the $M$ heads are chosen at random, and
  the variance of their forces about the full committee mean is rescaled by
  $M/\bigl(K(M-1)\bigr)$, which makes the estimate unbiased.
\item \emph{Randomised low-rank sketch with control variate.} The leading
  $r_0$ eigenvectors of the pooled design-split matrix
  $G=\sum_{x\in\mathrm{design}}A(x)$ are evaluated exactly, and the residual
  subspace orthogonal to those eigenvectors is sketched with the remaining
  $K-r_0$ random directions. This is the split between an exact leading
  subspace and a sketched residual that \citet{meyer2021hutchpp} show to be
  optimal for trace estimation. The estimate is unbiased and the random part
  sees only the energy left in the residual.
\item \emph{Deterministic leading subspace.} The $r_0$ leading eigenvectors of
  $G$ only, with no residual term: the estimate is
  $\sum_{k\le r_0}\lVert g_k\rVert^2/(M-1)$. The estimate is biased low by the
  energy in the residual subspace, involves no randomness, and costs $r_0+1$
  backward passes.
\end{enumerate}
Randomised estimators are run with ten fixed probe seeds, and the worst seed
is reported next to the mean, because a deployment draws one seed. Unless
stated otherwise, $K=4$ and $r_0=2$.

\subsubsection{The screening rule}

Let $S$ be the exact score, the global score unless stated otherwise, let
$\hat S$ be a cheap estimate of $S$, and let $\tau$ be a selection threshold
set on the design split as the 95th percentile of $S$ there, so that the exact
rule selects about 5\% of structures. The \emph{exact rule} computes $S(x)$ for
every structure and selects those with $S(x)\ge\tau$. The \emph{screening rule}
computes $\hat S(x)$ from $K$ projections; if $c_\alpha\hat S(x)<\tau$ the
structure is skipped, which means not selected and not evaluated further;
otherwise $S(x)$ is computed exactly and the structure is selected if and only
if $S(x)\ge\tau$.

The constant $c_\alpha$ is the split-conformal
one~\citep{vovk2005algorithmic,lei2018distribution}. With $n$ calibration
structures and ratios $\rho_i=S(x_i)/\hat S(x_i)$, $c_\alpha$ is the
$\lceil (n+1)(1-\alpha)\rceil$-th smallest $\rho_i$. If calibration and test
structures are exchangeable, then for a new structure
$\Pr\bigl[S(x)\le c_\alpha\hat S(x)\bigr]\ge 1-\alpha$, so the screening rule
skips a structure the exact rule would have selected with probability at most
$\alpha$. When $\lceil (n+1)(1-\alpha)\rceil>n$ no order statistic carries
this guarantee; the implementation then returns $c_\alpha=+\infty$, and the
rule skips nothing, which is the safe failure. The realised coverage of the
rule follows the $\mathrm{Beta}(k,\,n+1-k)$ law with
$k=\lceil (n+1)(1-\alpha)\rceil$. We verify that law by simulation rather
than assert it, at the smallest and largest calibration sets
($n=\fsCovNCalMin$ and $n=\fsCovNCalMax$) among the $\fsCovSplits$ registered
splits. The guarantee is weakest at the smallest, $n=\fsCovNCalMin$, where the
realised coverage is $\fsCovRealisedMin$ against the $\fsCovBetaMin$ the law
predicts. At the value $\alpha=0.05$ used throughout, $\fsCovInfeasible$ of
the $\fsCovSplits$ splits lack a feasible order statistic.

Two properties of the rule follow from the construction. First, a systematic
bias in $\hat S$ is absorbed by $c_\alpha$ exactly, because $c_\alpha$ is
fitted to the ratio $S/\hat S$; only the variance of $\hat S$ across
structures widens the ratio distribution and reduces skipping. Second, the set
selected by the screening rule is a subset of the set selected by the exact
rule, so the fidelity of the screening rule to the exact committee statistic
is measured as the drop in recall of the top-5\% error positives relative to
the exact rule. That drop is never negative. The benefit is measured as the
\emph{skip fraction}, the fraction of test structures for which the exact
computation was avoided.

What the rule costs when it does not skip must be stated, because this is the
obvious way to flatter such a method. There is no partial credit. A structure
the rule declines to skip is then evaluated exactly and in full, so the $K{+}1$
backward passes already spent on its cheap score are overhead rather than a
down payment on the exact computation. The skip fractions reported are
therefore counts of structures that avoided the exact computation altogether,
estimators are compared at matched numbers of backward passes, and the skip
fraction is reported in place of a speedup because the skip fraction is
independent of hardware, automatic-differentiation backend and batch size.
Wall-clock measurements are reported separately (Figure~\ref{fig:speedup}).

\subsection{Datasets, splits and contamination}
\label{sec:data}

\subsubsection{Datasets}

\emph{3BPA} is a 27-atom flexible organic molecule~\citep{kovacs2021beyond}.
The training set is the 100 frames sampled at 300~K that \citet{beck2025committee}
used for every 3BPA committee, with 300 further frames at 300~K for validation.
The test sets are the published frames at 300, 600 and 1200~K
(1669, 2138 and 2139 frames). Because all 3BPA committees were trained at
300~K, temperature is a graded distance from the training distribution with
the model, the molecule and $N$ held fixed.

\emph{rMD17} is the revised MD17 set of ten small organic
molecules~\citep{christensen2020role}. One committee was trained jointly on 50
frames of each of the ten molecules (500 frames; 100 frames per molecule for
validation). We score the held-out test frames of ethanol (9 atoms), aspirin
(21 atoms) and azobenzene (24 atoms), 1000 frames each. Every rMD17 result in
this paper therefore refers to one committee with disjoint heads, evaluated
on three molecules.

\emph{Liquid water} is a periodic box of 64 molecules (192 atoms, cubic cell
of edge 12.42~\AA). The committees were trained on 111 frames with 500 for
validation. Two test sets of 500 frames each at 300~K come from a classical
molecular dynamics (MD) trajectory and from a path-integral molecular dynamics
(PIMD) trajectory. The two differ only in nuclear quantum effects, which gives
a distribution shift with the temperature, the system and the cell held fixed.
All 3BPA and water frames and all checkpoints are the supporting data of
\citet{beck2025committee}; the rMD17 frames are those of
\citet{christensen2020role}. Every file is verified against a recorded SHA-256
hash before use.

\emph{MPtraj pool.} The materials pool is the set of $\fsMPtrajN$ periodic
structures from MPtraj~\citep{deng2023chgnet} that remained after the
query-by-committee selection of \citet{beck2025committee}, spanning 89
elements and 1 to 444 atoms with a median of 20 atoms. The pool contains many
near-equilibrium frames whose reference forces are close to zero. Structures of
different size cannot be stacked into one array, so the scan stores
$A(x)$ (Eq.~\ref{eq:A}) and the two maximum scores for each structure, and
the identity between the stored $A(x)$ and the direct computation of
Eq.~(\ref{eq:vd}) is asserted before the scan starts.

\subsubsection{Contamination between training and evaluation}

The MACE-MP-0b trunk was pretrained on MPtraj. The MPtraj pool is therefore
inside the training distribution of the trunk, and largely inside its training
set; the pool is out-of-sample only for the readout heads. Of the four
foundation committees, only QbC-8000 can be scored on the pool without
contamination of the heads, because the pool is by construction the complement
of the QbC selection: no structure of the QbC-8000 fine-tuning set appears in
the pool, which we checked by content hash. The random-8000 and max-force-8000
committees were fitted on structures most of which lie inside the pool
(\fsTODO{number of the 8,000 random-selected and of the 8,000
max-force-selected fine-tuning structures that lie inside the pool}), so those
committees are not scored on MPtraj. The submitted version of this paper did
not state the pretraining overlap; the out-of-distribution protocol of
Section~\ref{sec:protocols} exists because of that overlap. For 3BPA, rMD17
and water the test frames are disjoint from the training and validation frames
of every committee.

\subsubsection{Splits}

Three quantities are fitted in this paper: the leading subspace and the
threshold $\tau$ (on a design split), the constant $c_\alpha$ (on a calibration
split), and the exponent $\alpha$ of Eq.~(\ref{eq:salpha}) (on a design
split). Every number that involves a fitted quantity is reported on a test
split disjoint from the split the quantity was fitted on. Trajectory frames
are ordered in time, so a random split of frames can place near-duplicate
frames on both sides of a fit. Each trajectory is instead cut into ten
contiguous blocks, of which two serve as the design split, two as the
calibration split and six as the test split. The roles of the blocks are
permuted with a seed mixed with a hash of the system name, so that every role
samples the whole trajectory and a slow drift is not confounded with the role.
At each boundary where the role changes, a guard band of
$\lceil 2\tau_{\mathrm{int}}\rceil$ frames is dropped, where
$\tau_{\mathrm{int}}$ is the integrated autocorrelation time of the error
series measured by automatic windowing. The same $\tau_{\mathrm{int}}$ sets
the block length of the bootstrap in Section~\ref{sec:metric}. The MPtraj pool
is not a trajectory; there, a seeded random tenth of the pool is the design
split and the remainder is the test split. Each split is written to a manifest
whose content hash covers the structure identifiers and their roles, and each
manifest recomputes from its recorded seed, so a regenerated manifest hashes
identically if and only if the assignment is identical.

\subsection{Out-of-distribution and active-learning protocols}
\label{sec:protocols}

Two experiments were added for this revision. For each, the predictions, the
scoring script and the thresholds were written and registered by hash before
any structure was scored or any model trained.

\subsubsection{Out-of-distribution scan on OMat24}

The MPtraj result is open to two objections: the trunk was pretrained on the
pool, and the pool contains near-equilibrium frames with vanishing forces. The
scan addresses both. The evaluation data are three subsets of the validation
split of OMat24~\citep{barrosoluque2024omat24}, which post-dates the training
of MACE-MP-0 and which \citet{beck2025committee} used as their
out-of-distribution test: aimd-from-PBE-1000-nvt and aimd-from-PBE-1000-npt,
both ab initio molecular dynamics at 1000~K and far from equilibrium, and
rattled-1000-subsampled, rattled static frames. From each subset $10{,}000$
structures are drawn without replacement with a seeded generator (seed
20261111). Structures containing an element outside the 89-element table of
the model are dropped and counted. No structure in OMat24 was seen by the
trunk or by any readout head. All four foundation committees are scored. For
each committee and subset we report, for the top-5\% $e_{\max}$ positives,
the AUROC of the three committee scores and of the zero-cost baselines, and
the AUROC averaged within deciles of the reference-force magnitude, which asks
whether a committee score carries information beyond force magnitude. On the
in-distribution MPtraj pool the same scoring is repeated after excluding
structures whose largest reference-force norm is below 0.05~eV/\AA\ and
single-atom cells. Intervals are from a paired bootstrap with a recorded seed.
Five predictions were registered: (P1) the max-atom score reaches AUROC of at
least 0.70 on both molecular-dynamics subsets for the QbC-8000 committee;
(P2) on both molecular-dynamics subsets the predicted force norm is within
0.02 AUROC of, or above, the best committee score; (P3) on the filtered MPtraj
pool the margin of the predicted force norm over the best committee score
shrinks but stays positive; (P4) within reference-force deciles the max-atom
score exceeds AUROC 0.55 on both molecular-dynamics subsets; (P5) the
committee with the lowest out-of-distribution force RMSE has the highest
max-atom AUROC. Results are in Table~\ref{tab:ood}.

\subsubsection{One active-learning round on 3BPA}

Identifying high-error structures is not the same as acquiring the structures
that most improve a retrained model. One retraining round tests the link on
3BPA. The base training set is the 100 frames at 300~K used for every 3BPA
committee. The pool is the 2139 labelled frames of the 1200~K test set. Each
selection rule scores the pool with the shared-trunk committee with disjoint
heads and picks $k=50$ frames. The rules are: none (no frames added), the
global, max-component and max-atom scores, the predicted force norm, the
head-energy spread, and random. For each rule a single-head MACE model with
the independent-committee recipe of Section~\ref{sec:committees} (two layers,
32 channels, $L=1$, cutoff 6~\AA, energy and force weights 10 and 1000, EMA and
SWA, 625 epochs, float32, with the 300-frame validation file) is trained from
scratch on the 150 frames, with model seeds 0, 1 and 2; the random rule draws
a different subset for each seed. Each model is evaluated by force RMSE, force
MAE and per-atom energy MAE on (i) the 1820 pool frames that no rule selected,
a held-out set common to every rule, (ii) the 600~K test set and (iii) the
300~K test set. Five predictions were registered: (P1) every rule, random
included, lowers the held-out 1200~K force RMSE relative to the control with no
added frames; (P2) each of the global, max-component, max-atom and force-norm
rules gives a lower held-out force RMSE than random; (P3) the force-norm rule
is within 10\% (relative) of the best committee rule; (P4) at fixed $N$ the
global and max-component rules train models whose RMSE differ by less than the
across-seed standard deviation of either; (P5) on the 300~K test set no rule
improves the force RMSE over the control by more than 5\%. Results are in
Table~\ref{tab:al}.

\begin{table}[t]
\centering
\caption{Notation. Costs are counted in backward passes through a shared
trunk; the mean force is the one backward pass a simulation pays in any case.}
\label{tab:notation}
\small
\begin{tabular}{@{}llp{0.52\linewidth}@{}}
\toprule
Symbol & Name & Definition \\
\midrule
$N$ & atoms & number of atoms in a structure; $3N$ force coordinates \\
$M$ & heads & number of committee members ($M=8$ throughout) \\
$r=M-1$ & centred rank & dimension of head space after removing the mean \\
$F(x)$ & head forces & $3N\times M$ matrix, one column of forces per head \\
$\bar f(x)$ & committee mean & $M^{-1}\sum_m F_{\cdot m}(x)$; the force a simulation uses \\
$Q_c$, $A(x)$ & centred head space & Helmert basis of the centred subspace and the $r\times r$ matrix $Q_c^{\top}F^{\top}FQ_c$ (Eq.~\ref{eq:A}) \\
$v_d$ & head variance & unbiased variance over heads of coordinate $d$ (Eq.~\ref{eq:vd}) \\
$\sigma_a$ & per-atom std & mean over the three components of atom $a$ of $\sqrt{v_d}$ (Eq.~\ref{eq:sigma}) \\
$S_{\mathrm{glob}}$ & global score & $\sum_d v_d$; sum over all $3N$ coordinates \\
$S_{\mathrm{comp}}$ & max-component score & $\max_d\sqrt{v_d}$ \\
$S_{\mathrm{atom}}$ & max-atom score & $\max_a\sigma_a$ \\
$S_\alpha$ & size-normalised score & $S_{\mathrm{glob}}/(3N)^{\alpha}$ (Eq.~\ref{eq:salpha}) \\
$b_F$ & predicted force norm & $\max_a\lVert\bar f_a\rVert$; zero-cost baseline \\
$b_E$ & head-energy spread & standard deviation over heads of the total energy; zero-cost baseline \\
$\Delta$ & error & $\bar f-F_{\mathrm{ref}}$, error of the committee mean against the reference forces \\
$e_{\max}$ & largest per-atom error & $\max_a\lVert\Delta_a\rVert$ (primary target) \\
$e_{\mathrm{rmse}}$ & RMS error & root mean square of $\Delta_d$ over the $3N$ coordinates \\
$e_{\mathrm{maxcomp}}$ & largest component error & $\max_d\lvert\Delta_d\rvert$ \\
$e_{\mathrm{q95}}$ & 95th-percentile error & 95th percentile of $\lvert\Delta_d\rvert$ over the $3N$ coordinates \\
$g_k=Fw_k$ & projection & forces projected on head-space direction $w_k$; one backward pass each \\
$K$ & projections & number of head-space projections a cheap estimator uses; cost $K+1$ backward passes \\
$r_0$ & leading directions & number of exact leading eigendirections of the design-split matrix $G$ \\
$\tau$ & selection threshold & 95th percentile of the exact score on the design split \\
$c_\alpha$ & screening constant & split-conformal order statistic of $S/\hat S$ on the calibration split \\
$\alpha$ & miscoverage & allowed probability that a structure the exact rule selects is skipped ($0.05$) \\
\bottomrule
\end{tabular}
\end{table}
